% CP-VI-0016
% target_chapter: VI/14 — Sheafification
% source_unit_id: IMP-DL-0E3E7EC766-P016
% source: Downloads/theory-of-algebraic-geometry-1.tex L2904-L2955
% solution_provenance: SOURCE_EXPLICIT_SOLUTION

\subsection*{Exercise 4}

Let $\mathcal F$ be a presheaf of abelian groups on a topological space $X$.
Complete the proof of the existence of the associated sheaf $\mathcal F^+$
by showing:

\begin{enumerate}
	\item $\mathcal F^+$ is a sheaf.
	\item There is a natural morphism of presheaves
	\[
	\theta : \mathcal F \longrightarrow \mathcal F^+
	\]
	such that the pair $(\mathcal F^+,\theta)$ satisfies the desired universal property.
\end{enumerate}

Further, show that
\[
(\mathcal F^+)_p = \mathcal F_p
\qquad \text{for every } p \in X.
\]
Show that if
\[
f : \mathcal F \longrightarrow \mathcal G
\]
is a morphism of presheaves, then there is an induced morphism
\[
f^+ : \mathcal F^+ \longrightarrow \mathcal G^+
\]
with
\[
(f^+)_p = f_p.
\]

The universal property can be rephrased in terms of adjoint functors. Let
$\operatorname{PreSh}_X$ denote the category of presheaves of abelian groups on
$X$, and let $\operatorname{Sh}_X$ denote the category of sheaves of abelian
groups on $X$. There is an obvious inclusion functor
\[
i : \operatorname{Sh}_X \longrightarrow \operatorname{PreSh}_X.
\]
The above construction tells us that $i$ has an adjoint functor given by
\[
a(\mathcal F) = \mathcal F^+.
\]
This means that, for $\mathcal F$ a presheaf and $\mathcal G$ a sheaf, there
are natural bijections
\[
\operatorname{Hom}_{\operatorname{Sh}_X}(\mathcal F^+,\mathcal G)
\cong
\operatorname{Hom}_{\operatorname{PreSh}_X}(\mathcal F,i\mathcal G).
\]

\subsection*{Solution}

We construct the sheafification of $\mathcal F$ by using stalks.

For each point $p \in X$, let
\[
\mathcal F_p
=
\varinjlim_{p \in U} \mathcal F(U)
\]
be the stalk of $\mathcal F$ at $p$.

Elements of $\mathcal F_p$ are germs. If $s \in \mathcal F(U)$ and
$p \in U$, we denote its germ at $p$ by
\[
s_p \in \mathcal F_p.
\]

Define the disjoint union
\[
E_{\mathcal F}
:=
\coprod_{p \in X} \mathcal F_p.
\]
There is a natural projection
\[
\pi : E_{\mathcal F} \longrightarrow X,
\qquad
\pi(a) = p
\quad \text{if } a \in \mathcal F_p.
\]

For every open set $U \subseteq X$ and every section $s \in \mathcal F(U)$,
define a map
\[
\widetilde{s} : U \longrightarrow E_{\mathcal F}
\]
by
\[
\widetilde{s}(p) = s_p.
\]

We topologize $E_{\mathcal F}$ by declaring the sets
\[
\widetilde{s}(U)
=
\{s_p : p \in U\}
\]
to be a basis of open sets.

Then define
\[
\mathcal F^+(U)
\]
to be the set of all continuous sections of $\pi$ over $U$, i.e.
\[
\mathcal F^+(U)
=
\{\sigma : U \to E_{\mathcal F}
\mid \pi \circ \sigma = \operatorname{id}_U,\ \sigma
\text{ is continuous}\}.
\]

For an inclusion $V \subseteq U$, the restriction map is ordinary restriction:
\[
\rho^U_V : \mathcal F^+(U) \longrightarrow \mathcal F^+(V),
\qquad
\sigma \longmapsto \sigma|_V.
\]

Thus $\mathcal F^+$ is at least a presheaf.

\subsection*{1. The presheaf $\mathcal F^+$ is a sheaf}

Let $U \subseteq X$ be open and let $\{U_i\}_{i \in I}$ be an open cover of $U$.

Suppose we have sections
\[
\sigma_i \in \mathcal F^+(U_i)
\]
such that
\[
\sigma_i|_{U_i \cap U_j}
=
\sigma_j|_{U_i \cap U_j}
\]
for all $i,j$.

We define a map
\[
\sigma : U \longrightarrow E_{\mathcal F}
\]
by setting
\[
\sigma(p) = \sigma_i(p)
\quad \text{if } p \in U_i.
\]

This is well-defined because if $p \in U_i \cap U_j$, then
\[
\sigma_i(p) = \sigma_j(p).
\]

Also
\[
\pi(\sigma(p)) = p,
\]
because each $\sigma_i$ is a section of $\pi$.

It remains to show that $\sigma$ is continuous.

Continuity is local on the domain. Since each $\sigma_i$ is continuous and
\[
\sigma|_{U_i} = \sigma_i,
\]
it follows that $\sigma$ is continuous.

Hence
\[
\sigma \in \mathcal F^+(U).
\]

Moreover, $\sigma$ is the unique section with
\[
\sigma|_{U_i} = \sigma_i
\]
for every $i$, because the $U_i$ cover $U$.

Therefore $\mathcal F^+$ satisfies the sheaf gluing axiom.

The locality axiom is also immediate: if
\[
\sigma,\tau \in \mathcal F^+(U)
\]
and
\[
\sigma|_{U_i} = \tau|_{U_i}
\]
for every $i$, then for every $p \in U$ there exists $i$ with $p \in U_i$,
so
\[
\sigma(p) = \tau(p).
\]
Thus $\sigma=\tau$.

Therefore $\mathcal F^+$ is a sheaf.

\subsection*{2. The natural morphism $\theta : \mathcal F \to \mathcal F^+$}

For every open set $U \subseteq X$, define
\[
\theta_U : \mathcal F(U) \longrightarrow \mathcal F^+(U)
\]
by
\[
\theta_U(s) = \widetilde{s},
\]
where
\[
\widetilde{s}(p) = s_p.
\]

This is well-defined because $\widetilde{s}$ is continuous by construction of
the topology on $E_{\mathcal F}$.

The maps $\theta_U$ are compatible with restrictions. Indeed, if
$V \subseteq U$, then for $s \in \mathcal F(U)$ and $p \in V$,
\[
\theta_V(s|_V)(p)
=
(s|_V)_p
=
s_p
=
\theta_U(s)(p).
\]
Hence
\[
\theta_V(s|_V)
=
\theta_U(s)|_V.
\]

Thus $\theta$ is a morphism of presheaves:
\[
\theta : \mathcal F \longrightarrow \mathcal F^+.
\]

\subsection*{3. The stalks of $\mathcal F^+$}

We now show that
\[
(\mathcal F^+)_p \cong \mathcal F_p.
\]

Define
\[
\Phi_p : (\mathcal F^+)_p \longrightarrow \mathcal F_p
\]
by
\[
\Phi_p(\sigma_p) = \sigma(p).
\]

This makes sense because if
\[
\sigma \in \mathcal F^+(U)
\]
and $p \in U$, then
\[
\sigma(p) \in \mathcal F_p.
\]

We show that $\Phi_p$ is an isomorphism.

\subsubsection*{Surjectivity}

Let
\[
a \in \mathcal F_p.
\]
Then $a$ is represented by some section
\[
s \in \mathcal F(U)
\]
with $p \in U$, so
\[
a = s_p.
\]

Consider
\[
\widetilde{s} \in \mathcal F^+(U).
\]
Then
\[
\Phi_p((\widetilde{s})_p)
=
\widetilde{s}(p)
=
s_p
=
a.
\]

Thus $\Phi_p$ is surjective.

\subsubsection*{Injectivity}

Suppose
\[
\sigma,\tau \in \mathcal F^+(U)
\]
and suppose their germs at $p$ satisfy
\[
\Phi_p(\sigma_p) = \Phi_p(\tau_p).
\]
That means
\[
\sigma(p) = \tau(p)
\quad \text{in } \mathcal F_p.
\]

Because $\sigma$ and $\tau$ are continuous sections of the espace étalé,
there exist neighborhoods $V,W \subseteq U$ of $p$ and presheaf sections
\[
s \in \mathcal F(V),
\qquad
t \in \mathcal F(W),
\]
such that
\[
\sigma|_V = \widetilde{s},
\qquad
\tau|_W = \widetilde{t}.
\]

The equality
\[
\sigma(p)=\tau(p)
\]
means
\[
s_p = t_p
\quad \text{in } \mathcal F_p.
\]
Therefore there exists an open neighborhood
\[
N \subseteq V \cap W
\]
of $p$ such that
\[
s|_N = t|_N.
\]

Consequently
\[
\sigma|_N = \widetilde{s}|_N
=
\widetilde{t}|_N
=
\tau|_N.
\]

Thus $\sigma_p=\tau_p$ in $(\mathcal F^+)_p$.

Therefore $\Phi_p$ is injective.

Hence
\[
(\mathcal F^+)_p \cong \mathcal F_p.
\]

Under this identification,
\[
\theta_p : \mathcal F_p \longrightarrow (\mathcal F^+)_p
\]
is the identity map.

So one usually writes
\[
(\mathcal F^+)_p = \mathcal F_p.
\]

\subsection*{4. Universal property}

Let $\mathcal G$ be a sheaf and let
\[
\varphi : \mathcal F \longrightarrow \mathcal G
\]
be a morphism of presheaves.

We must construct a unique morphism of sheaves
\[
\varphi^+ : \mathcal F^+ \longrightarrow \mathcal G
\]
such that
\[
\varphi^+ \circ \theta = \varphi.
\]

Let $U \subseteq X$ be open and let
\[
\sigma \in \mathcal F^+(U).
\]

For each point $p \in U$, the value $\sigma(p)$ lies in $\mathcal F_p$.
The morphism $\varphi$ induces a morphism on stalks
\[
\varphi_p : \mathcal F_p \longrightarrow \mathcal G_p.
\]

We would like to define a section of $\mathcal G$ whose germ at $p$ is
\[
\varphi_p(\sigma(p)).
\]

This is possible because $\mathcal G$ is a sheaf.

Indeed, since $\sigma$ is continuous, for every $p \in U$ there is an open
neighborhood $U_p \subseteq U$ and a section
\[
s_p \in \mathcal F(U_p)
\]
such that
\[
\sigma|_{U_p} = \widetilde{s_p}.
\]

Then define locally
\[
g_p := \varphi_{U_p}(s_p) \in \mathcal G(U_p).
\]

On overlaps $U_p \cap U_q$, the local sections $g_p$ and $g_q$ have the same
germs at every point. Since $\mathcal G$ is a sheaf, equality of germs at every
point implies equality locally, and therefore the $g_p$ glue uniquely to a
section
\[
g \in \mathcal G(U).
\]

Define
\[
\varphi^+_U(\sigma) := g.
\]

Equivalently, $\varphi^+_U(\sigma)$ is the unique section of $\mathcal G(U)$
satisfying
\[
\bigl(\varphi^+_U(\sigma)\bigr)_p
=
\varphi_p(\sigma(p))
\qquad
\text{for all } p \in U.
\]

This construction is compatible with restrictions, so it gives a morphism of
sheaves
\[
\varphi^+ : \mathcal F^+ \longrightarrow \mathcal G.
\]

Now let $s \in \mathcal F(U)$. Then
\[
\theta_U(s)=\widetilde{s}.
\]
For every $p \in U$,
\[
\bigl(\varphi^+_U(\widetilde{s})\bigr)_p
=
\varphi_p(s_p)
=
(\varphi_U(s))_p.
\]
Since $\mathcal G$ is a sheaf, sections with the same germs at all points are
equal. Hence
\[
\varphi^+_U(\theta_U(s))
=
\varphi_U(s).
\]
Therefore
\[
\varphi^+ \circ \theta = \varphi.
\]

Finally, uniqueness follows because the local sections $\widetilde{s}$ generate
$\mathcal F^+$ locally. More explicitly, every section of $\mathcal F^+$ is
locally of the form $\widetilde{s}$. Therefore any morphism
\[
\psi : \mathcal F^+ \to \mathcal G
\]
satisfying
\[
\psi \circ \theta = \varphi
\]
must agree with $\varphi^+$ locally, hence globally by the sheaf axiom.

Thus $(\mathcal F^+,\theta)$ satisfies the universal property of
sheafification.

\subsection*{5. Functoriality}

Let
\[
f : \mathcal F \longrightarrow \mathcal G
\]
be a morphism of presheaves.

Applying the universal property to the morphism
\[
\mathcal F \xrightarrow{f} \mathcal G
\xrightarrow{\theta_{\mathcal G}} \mathcal G^+
\]
gives a unique morphism of sheaves
\[
f^+ : \mathcal F^+ \longrightarrow \mathcal G^+
\]
such that
\[
f^+ \circ \theta_{\mathcal F}
=
\theta_{\mathcal G} \circ f.
\]

On stalks this gives
\[
(f^+)_p = f_p.
\]

Indeed, under the identifications
\[
(\mathcal F^+)_p \cong \mathcal F_p,
\qquad
(\mathcal G^+)_p \cong \mathcal G_p,
\]
the induced map is exactly the map on germs
\[
f_p : \mathcal F_p \longrightarrow \mathcal G_p.
\]

Therefore sheafification is functorial:
\[
\mathcal F \longmapsto \mathcal F^+,
\qquad
f \longmapsto f^+.
\]

\subsection*{6. Adjointness}

Let
\[
i : \operatorname{Sh}_X \longrightarrow \operatorname{PreSh}_X
\]
be the inclusion functor.

The universal property says exactly that for every presheaf $\mathcal F$ and
every sheaf $\mathcal G$, composition with
\[
\theta : \mathcal F \to i(\mathcal F^+)
\]
induces a bijection
\[
\operatorname{Hom}_{\operatorname{Sh}_X}(\mathcal F^+,\mathcal G)
\cong
\operatorname{Hom}_{\operatorname{PreSh}_X}(\mathcal F,i\mathcal G).
\]

Hence the sheafification functor
\[
a : \operatorname{PreSh}_X \longrightarrow \operatorname{Sh}_X,
\qquad
a(\mathcal F)=\mathcal F^+
\]
is left adjoint to the inclusion functor:
\[
a \dashv i.
\]

Thus sheafification is the best sheaf approximation to a presheaf.

\[
\boxed{
	\operatorname{Hom}_{\operatorname{Sh}_X}(\mathcal F^+,\mathcal G)
	\cong
	\operatorname{Hom}_{\operatorname{PreSh}_X}(\mathcal F,\mathcal G)
}
\]
whenever $\mathcal G$ is already a sheaf.

\newpage
